\section{Game Routes}
\label{sec:cses-game-routes}

\problemstatement{Given a \emph{directed acyclic} graph of $n$ levels and
$m$ edges, count the number of distinct routes from level $1$ to level $n$.
Output the count modulo $10^9+7$.}

\begin{patternbox}
Counting paths in a \textbf{DAG} is a topological-order DP: the number of
routes to a node is the sum of the routes to its predecessors. Acyclicity
guarantees each predecessor's count is final before it is used.
\end{patternbox}

\subsection*{Sum over predecessors}

Let $\textit{ways}[v]$ be the number of routes from $1$ to $v$, with
$\textit{ways}[1]=1$. In topological order, push each node's count forward:
for every edge $u\to v$, add $\textit{ways}[u]$ to $\textit{ways}[v]$
(modulo $10^9+7$). Since $u$ comes before $v$ topologically,
$\textit{ways}[u]$ is complete when added. The answer is $\textit{ways}[n]$.

\begin{ideabox}[Every Route Has a Unique Last Edge]
Classifying routes to $v$ by their final edge $u\to v$ partitions them
without overlap, so summing $\textit{ways}[u]$ over all predecessors $u$
counts each route exactly once -- the DAG analogue of the coin-counting
recurrence.
\end{ideabox}

\subsection*{Algorithm}

\begin{enumerate}
  \item Topologically sort. $\textit{ways}[1]=1$.
  \item In topological order, for each edge $u\to v$ add $\textit{ways}[u]$
        to $\textit{ways}[v]$ modulo $10^9+7$.
  \item Output $\textit{ways}[n]$.
\end{enumerate}

\subsection*{Dry run}

\dryrun{$1\!\to\!2$, $1\!\to\!3$, $2\!\to\!4$, $3\!\to\!4$; $n=4$.}

\begin{itemize}
  \item $\textit{ways}[1]=1$, $[2]=1$, $[3]=1$, $[4]=\textit{ways}[2]+
        \textit{ways}[3]=2$.
  \item Two routes: $1\to2\to4$ and $1\to3\to4$.
\end{itemize}

\subsection*{C++ implementation}

\begin{lstlisting}
#include <bits/stdc++.h>
using namespace std;
const long long MOD = 1e9 + 7;

int main() {
    int n, m;
    cin >> n >> m;
    vector<vector<int>> adj(n + 1);
    vector<int> indeg(n + 1, 0);
    for (int i = 0; i < m; i++) {
        int a, b; cin >> a >> b;
        adj[a].push_back(b);
        indeg[b]++;
    }

    vector<int> order, deg = indeg;
    queue<int> q;
    for (int i = 1; i <= n; i++) if (deg[i] == 0) q.push(i);
    while (!q.empty()) {
        int u = q.front(); q.pop(); order.push_back(u);
        for (int v : adj[u]) if (--deg[v] == 0) q.push(v);
    }

    vector<long long> ways(n + 1, 0);
    ways[1] = 1;
    for (int u : order)
        for (int v : adj[u])
            ways[v] = (ways[v] + ways[u]) % MOD;

    cout << ways[n] << "\n";
    return 0;
}
\end{lstlisting}

\begin{complexitybox}
\textbf{Time Complexity.} $O(n+m)$. \textbf{Space Complexity.} $O(n+m)$.
\end{complexitybox}

\begin{takeawaybox}
DAG path counting sums predecessor counts in topological order. Longest
Flight Route and Game Routes are the same sweep with $\max{+}1$ versus
sum -- the recurring pairing of optimisation and counting on the same DP
skeleton.
\end{takeawaybox}
